\chapter{Introduzione}
\label{cap:rad-introduzione}

\section{Scopo del documento}

Questo documento definisce i requisiti della piattaforma Tech4All. Ne descrive
il dominio applicativo, gli attori che vi interagiscono, le funzionalità
offerte e le qualità che il sistema deve esibire, senza entrare nel merito
delle scelte realizzative, che sono trattate nel System Design Document.

Il documento segue l'impostazione di Bruegge e Dutoit~\cite{bruegge2010}: dalla
descrizione del sistema esistente si passa al sistema proposto, dai requisiti
si risale agli scenari e da questi ai modelli di analisi. L'unico scostamento
dalla struttura del testo di riferimento è la promozione dei modelli del
sistema a capitolo autonomo (\cref{cap:rad-modelli}), scelta dettata dalla
leggibilità e priva di conseguenze metodologiche.

Il perimetro del sistema è definito in modo esplicito nelle due direzioni:
\cref{sec:rad-fuori-perimetro} elenca le funzionalità che il progetto ha
deciso di non realizzare, con la motivazione di ciascuna esclusione. Un
requisito assente per scelta e un requisito dimenticato producono lo stesso
sistema, ma non la stessa possibilità di discuterne.

\section{Ambito del sistema}

Tech4All è un'applicazione web per l'alfabetizzazione digitale. Si rivolge a
persone con poca dimestichezza con gli strumenti informatici e propone un
percorso in tre momenti: la lettura di un \term{tutorial}, la verifica
tramite un \term{quiz} a scelta multipla e il riconoscimento dei progressi
attraverso \term{badge} assegnati automaticamente.

Rientrano nell'ambito del sistema:

\begin{itemize}
  \item la pubblicazione e la consultazione di un catalogo di tutorial
        classificati per categoria;
  \item la verifica dell'apprendimento tramite quiz corretti dal sistema;
  \item la raccolta di valutazioni degli utenti sui contenuti;
  \item il riconoscimento dei traguardi raggiunti;
  \item l'amministrazione di contenuti, obiettivi e account.
\end{itemize}

Non rientrano nell'ambito del sistema la produzione dei contenuti didattici
(redatti dagli amministratori con gli strumenti offerti dalla piattaforma),
l'erogazione di percorsi certificati e qualunque forma di tutoraggio umano.

\section{Obiettivi e criteri di successo}
\label{sec:rad-obiettivi}

Gli obiettivi del progetto sono formulati in modo da poter essere verificati.
A ciascuno corrisponde un criterio osservabile, riportato in
\cref{tab:rad-obiettivi}.

\begin{table}[H]
\centering
\begin{tabular}{@{}l L{5.6cm} L{6.3cm}@{}}
\toprule
\textbf{Id} & \textbf{Obiettivo} & \textbf{Criterio di successo} \\
\midrule
\id{OB\_1} &
Rendere i contenuti accessibili senza barriere d'ingresso &
Il catalogo e il singolo tutorial sono consultabili senza autenticazione. \\

\id{OB\_2} &
Permettere all'utente di verificare ciò che ha appreso &
Ogni tutorial può avere un quiz; l'esito è calcolato dal sistema e
comunicato immediatamente. \\

\id{OB\_3} &
Restituire all'utente una misura del proprio progresso &
Il numero di quiz superati e i badge conseguiti sono consultabili
nell'area personale. \\

\id{OB\_4} &
Raccogliere il giudizio degli utenti sui contenuti &
Ogni utente può lasciare al più un feedback per tutorial; la media è
mostrata nel catalogo. \\

\id{OB\_5} &
Garantire che i dati personali siano trattati in modo sicuro &
Le password sono conservate solo come hash; nessuna operazione di
modifica è possibile senza una sessione valida. \\

\id{OB\_6} &
Impedire che il quiz possa essere superato aggirando la verifica &
Le soluzioni non sono trasmesse al client prima della consegna e la
correzione avviene solo sul server. \\
\bottomrule
\end{tabular}
\caption{Obiettivi del sistema e relativi criteri di successo.}
\label{tab:rad-obiettivi}
\end{table}

\section{Struttura del documento}

Il \cref{cap:rad-sistema-corrente} descrive la situazione che il sistema
intende migliorare. Il \cref{cap:rad-sistema-proposto} introduce il sistema
proposto, ne individua gli attori ed elenca i requisiti funzionali e non
funzionali. Il \cref{cap:rad-modelli} raccoglie i modelli di analisi:
scenari, casi d'uso, modello a oggetti, modello dinamico e interfaccia
utente. Il \cref{cap:glossario} fissa il linguaggio del dominio.

Le due appendici raccolgono il materiale di consultazione:
l'appendice~\ref{cap:rad-app-requisiti} riporta i requisiti in forma tabellare
completa, l'appendice~\ref{cap:rad-app-tracciabilita} la matrice che lega requisiti,
casi d'uso e scenari.

\section{Riferimenti}

Il documento fa riferimento a Bruegge e Dutoit~\cite{bruegge2010} per il
processo e la notazione, a Fowler~\cite{fowler2003} per l'uso di UML e a
ISO/IEC 25010~\cite{iso25010} per la classificazione delle qualità del
software, adottata come complemento al modello FURPS+.

Gli altri documenti del progetto sono il System Design Document, che descrive
architettura e progettazione degli oggetti, e il Test Document, che descrive
strategia, casi di test e risultati.
